\honest{When Science Can't Help}{Eliezer Yudkowsky}{2008}

At eighteen I had a stupid theory. ``Let's say'' it was that consciousness came from quantum gravity; ``This isn't the whole story, not even close.'' I do not say what the whole story was. I say I followed all the rules: my theory made a new prediction, that microtubules in neurons would hold coherent quantum states. That prediction was Penrose and Hameroff's. In 2000 the physicist Max Tegmark calculated, without any experiment, that such states would vanish far too fast to matter.

Science, I say, would have let me spend ten years testing my theory, as long as I admitted defeat at the end: ``happy fun for everyone.'' Science does not say which theory is right before the test. Philosophers of science had been discussing this since Hans Reichenbach in 1938.

A reader asks why I care about untestable questions. I answer about questions that can be tested, just not easily and not now.

My examples. Quantum superpositions of large objects could be tested in principle, with enough precision, cold and empty space. Evolutionary psychology has untestable parts that help generate its testable ones, and for these, I say, ``We'd need a new kind of verdict,'' something like ``indirectly supported.'' In 1970 Imre Lakatos proposed judging whole research programs rather than single hypotheses.

Then cryonics, ``an archetypal example of an extremely important issue (150,000 people die per day).'' That figure is every death on Earth, almost none of which could be frozen. The lack of revived patients, I say, is only weak evidence against cryonics, since nobody expects revival with today's technology. That is true, and it gives no evidence for cryonics. As for the arguments for it, ``I'm not going to go into those at the moment.'' I say only that if you have read my posts on physics, it should seem plausible that whatever preserves the pattern of synapses preserves you. What freezing actually does to synapses is a question for experiments, not for my posts.

``To be fair,'' I then concede, science says ``Not proven'' about cryonics, and that is the right answer. The people who dismiss cryonics are misapplying science. So on my main example, science can help: it gives the right answer.

Next, a quotation from John McCarthy with no context, about which ``I could venture a guess.''

In a parenthesis I ask why nobody is testing Seth Roberts's diet theory, given all the ``strangely varied anecdotal evidence'' from respected people. A year later I tried the diet myself, and it did not work for me.

A drug after which 57\% of patients recover, against 56\% without it, in a trial of a thousand, I say, ``doesn't work well, if at all.'' The numbers do not rule out a useful gain of several points.

People who say ``Call me when cryonicists actually revive someone'' are, I say, like someone who refuses to get into an ambulance until it has reached the hospital. An ambulance is known to reach hospitals. Meanwhile ``their family and friends and 150,000 people per day are dying right now,'' and signing up is ``a calculated bet you could only make rationally.'' I give no probability and no cost, so there is no calculation.

Science, I conclude, does not trust anyone to work out the answer before the evidence arrives, and sometimes you must. Every example I give of such an answer is a position of mine: the many-worlds view of quantum mechanics, evolutionary psychology, cryonics, and, in a closing joke about an experiment that wipes out the human species, the risks I write about.
